\setcounter{deso}{4}
\begin{name}
	{\tenchude}
	{\tendethi}
	{\tentruong}
	{\thoigian}
\end{name}
\begin{bt}[1,5 điểm]~
\begin{enumerate}[1.]
    \item Tính giá trị biểu thức $A=\sqrt{(3-\sqrt{7})^2}-\sqrt{28}+\sqrt{3}(\sqrt{27}+\sqrt{12})$.
    \item Giải phương trình $\dfrac{1}{2}x-1=-\dfrac{2}{3}(2x-6)$.
    \item Giải bất phương trình $2(x-1) \ge 3-2x$.
\end{enumerate}
\end{bt}

\begin{bt}[2,0 điểm]~
\begin{enumerate}[1.]
    \item Cho hàm số $y=ax^2$ với $a \ne 0$ có đồ thị cắt đường thẳng $y=x$ tại hai điểm có hoành độ là $0$ và $1$. Xác định hệ số $a$.
    \item Cho biểu thức $M=\dfrac{x+2\sqrt{x}+1}{x\sqrt{x}+1} : \dfrac{x-1}{x-\sqrt{x}+1}$
    \begin{enumerate}
        \item Rút gọn biểu thức $M$.
        \item Tìm tất cả giá trị của $x$ để $M$ là số nguyên.
    \end{enumerate}
\end{enumerate}
\end{bt}
\begin{bt}[1,5 điểm]%[Đề Thi Tuyển Sinh 10, Bà Rịa - Vũng Tàu 2025-2026]%[DethiEX-9CT-2025-2026, Đỗ Nam]%[9D5H1-1]%[9D6H2-2]	
 	Một người đi xe máy từ $A$ đến $B$ cách nhau $100$\,km. Lúc từ $B$ trở về $A$, người đó đi với tốc độ nhanh hơn lúc đi là $10$\,km/h. Biết tổng thời gian cả đi và về là $4$ giờ $30$ phút. Tính tốc độ của xe máy lúc đi.
 \loigiai{
 Gọi tốc độ của xe máy lúc đi là $x$\,(km/h), $(x>0)$.\\
	Tốc độ của xe máy lúc về là $x+10$\,(km/h).
	Thời gian xe máy đi từ $A$ đến $B$ là $\dfrac{100}{x}$\,(h).
	Thời gian xe máy đi từ $B$ về $A$ là $\dfrac{100}{x+10}$\,(h).\\
	Đổi $4$ giờ $30$ phút $=\dfrac{9}{2}$ giờ.\\
	Vì tổng thời gian cả đi lẫn về là $4$ giờ $30$ phút nên ta có phương trình sau
	\begin{eqnarray*}
		&&\dfrac{100}{x}+\dfrac{100}{x+10}=\dfrac{9}{2}\\
		&&\dfrac{100\cdot2\cdot(x+10)}{2x(x+10)}+\dfrac{100\cdot2\cdot x}{2x(x+10)}=\dfrac{9x(x+10)}{2x(x+10)}\\
		&&100\cdot2\cdot(x+10)+100\cdot2\cdot x=9x(x+10)\\
		&&200x+2\,000+200x=9x^2+90x\\
		&&9x^2+90x-200x-2\,000-200x=0\\
		&&9x^2-310x-2\,000=0\quad (*)
	\end{eqnarray*}
	Ta có $\Delta=(310)^2-4\cdot9\cdot(-2\,000)=168\,100$.\\
	Vì $\Delta>0$ nên phương trình $(*)$ có hai nghiệm phân biệt là
	\[\hoac{&x_1=\dfrac{-b+\sqrt{\Delta}}{2a}=\dfrac{-(-310)+\sqrt{168\,100}}{2\cdot9}=40\text{ (thỏa mãn điều kiện)}\\&x_2=\dfrac{-b-\sqrt{\Delta}}{2a}=\dfrac{-(-310)-\sqrt{168\,100}}{2\cdot9}=-\dfrac{50}{9}\text{ (loại).}}\]
	Vậy tốc độ của xe máy lúc đi là $40$\,km/h.
}
\end{bt} 


\begin{bt}[1,0 điểm]%[9CT-2025-2026-TS10-BenTre-Nguyễn Mạnh Hà]%[9D5H1-1]%[9D6H2-2]%
 Một hộp có chứa $5$ viên bi cùng loại, trong đó có hai viên bi màu vàng lần lượt ghi các số $1$; $2$ và ba viên bi màu đỏ lần lượt ghi các số $3$; $4$; $5$. Lấy ra ngẫu nhiên đồng thời hai viên bi từ hộp. Tính xác suất của biến cố $A\colon$\lq\lq Hai viên bi được lấy ra khác màu\rq\rq.

	\loigiai{%
		Không gian mẫu: $\Omega=\left\{(1,2);(1,3);(1,4);(1,5);(2,3);(2,4);(2,5);(3,4);(3,5);(4,5)\right\}$.\\
			Số phần tử của không gian mẫu là $n\left(\Omega\right)=10$.\\
			Ta có $A=\left\{(1,3);(1,4);(1,5);(2,3);(2,4);(2,5)\right\}$, suy ra $n(A)=6$.\\
			Tính xác suất của biến cố $A$ là $P(A)=\dfrac{n(A)}{n(\Omega)}=\dfrac{6}{10}=\dfrac{3}{5}$.
	}
\end{bt}
%%%==============HetBai_BT6==============%%%

%%%==============Bai_BT7==============%%%
\begin{bt}[1,0 điểm]%[Đề tham khảo tuyển sinh vào 10 - Quận 8 - Đề số 1]%[TS10K6]
Từ một khúc gỗ hình trụ cao 15 cm, người ta tiện thành một hình nón có đáy là hình tròn bằng với đáy hình trụ, chiều cao của hình nón bằng chiều cao của hình trụ. Biết phần gỗ bỏ đi có thể tích là 3610$\pi$ (cho biết $\pi \approx 3{,}14$. Công thức tính thể tích hình trụ : $V = \pi R^2h$, thể tích hình nón: $V = \dfrac 13 \pi R^2h$ với $R$ là bán kính đáy, $h$ là chiều cao khúc gỗ). Tính thể tích khúc gỗ hình trụ, (làm tròn tới hàng đơn vị).
\dapso{$V\approx 17003$ (m$^3$).}
\loigiai{
\immini
{
Phần gỗ bỏ đi có thể tích là $3610\pi$ suy ra\\
$\pi R^2h - \dfrac{1}{3} \pi R^2h=3610\pi\Leftrightarrow\dfrac{2}{3}\cdot R^2h=3610\\\Leftrightarrow R^2=\dfrac{3610\cdot3}{2h}=\dfrac{3610\cdot 3}{2\cdot 15}=361\Rightarrow R=19$ (cm).\\
Vậy thể tích khúc gỗ hình trụ là $V=\pi\cdot 19^2\cdot 15\approx 17003$ (m$^3$).
}
{
\begin{tikzpicture}[scale=.7, font=\footnotesize, line join=round, line cap=round, >=stealth]
\def\x{2.5} % Bán kính trụ lớn
\pgfmathsetmacro{\y}{\x/4} % Bán kính trục bé
\def\h{5} % Chiều cao
\coordinate[label=right:$A$] (A) at (\x,0);
\coordinate[label=left:$B$] (B) at (-\x,0);
\coordinate[label=below left:$O$] (O) at (0,0);
\coordinate (A') at ($(A)+(90:\h)$);
\coordinate (A1) at ($(A)+(1,0)$);
\coordinate (A2) at ($(A')+(1,0)$);
\coordinate[label=left:$O'$] (O') at ($(O)+(90:\h)$);
\draw (A) arc (0:-180:{\x} and {\y})--($(A')!2!(O')$) arc (180:0:{\x} and {\y}) arc (0:-180:{\x} and {\y}) (A)--(A');
\draw[dashed] (O')--(O)--(A) arc (0:180:{\x} and {\y}) (A)--(O')--(B)--(O);
\draw[<->] (A1)--(A2)node[pos=.25,rotate=-90,above left]{\scriptsize$h=15$\,cm};
\foreach \diem in {A,B,O,O'}\fill (\diem)circle(1.5pt);
\end{tikzpicture}
}
}
\end{bt}

\begin{bt}[3,0 điểm]%[Đề Thi Tuyển Sinh 10, Bà Rịa - Vũng Tàu 2025-2026]%[DethiEX-9CT-2025-2026, Đỗ Nam]%[9H3V4-6]	
	Cho tam giác nhọn $ABC$ $(AB<AC)$ nội tiếp đường tròn $(O)$ và có đường cao $AD$, $BE$, $CF$ cắt nhau tại $H$.
	\begin{enumerate}
		\item Chứng minh tứ giác $BEFC$ nội tiếp.
		\item Vẽ đường kính $AT$ của đường tròn $(O)$. Chứng minh $\triangle ADB$ đồng dạng $\triangle ACT$ và $2\widehat{HEF}+\widehat{AOC}=180^{\circ}$.
		\item Vẽ $CI$ vuông góc với $AT$ tại $I$. Gọi $M$ là trung điểm của $BC$. Chứng minh ba điểm $F$, $M$, $I$ thẳng hàng.
	\end{enumerate}
	\loigiai{
		\begin{center}
			\begin{tikzpicture}[scale=1, font=\footnotesize, line join=round, line cap=round, >=stealth]
				\def\r{5}
				\path
				(0,0) coordinate (O)
				($(O)+(-30:\r)$) coordinate (C)
				($(O)+(-150:\r)$) coordinate (B)
				($(O)+(110:\r)$) coordinate (A)
				($(O)+(-70:\r)$) coordinate (T)
				($(A)!(B)!(C)$) coordinate (E)
				($(A)!(C)!(B)$) coordinate (F)
				($(C)!(A)!(B)$) coordinate (D)
				($(T)!(C)!(A)$) coordinate (I)
				%($(M)!0.5!(B)$) coordinate (P)
				(intersection of C--F and A--D) coordinate (H)
				(intersection of F--I and C--B) coordinate (M)
				%(intersection of P--O and D--M) coordinate (E)
				;
				\draw
				(O) circle (\r)
				(A)--(B)--(C)--cycle
				(T)--(B)--(E)--(F)--(C)--(I)--(F) (M)--(O)--(C)--(T)--(A)--(D)
				;
				\foreach \a/\b/\c in {A/F/C,C/E/B,A/D/B,C/I/A}{
					\draw pic[draw, angle radius=2mm,fill=gray!50]{right angle=\a--\b--\c};
				}
				\foreach \x/\g in {A/110,B/-135,C/-45,E/30,F/150,O/30,D/-90,H/125,T/-60,I/-150,M/-90} \draw[fill=black] (\x) circle (1pt) +(\g:3mm) node {$\x$};
			\end{tikzpicture}
		\end{center}
		\begin{enumerate}
			\item Do $BE$, $CF$ là đường cao $\triangle ABC$ nên $\triangle BEC$ vuông tại $E$ và $\triangle BFC$ vuông tại $F$.\\
			Ta có 
			\begin{itemize}
				\item $\triangle BEC$ vuông tại $E$ nên $B$, $E$, $C$ cùng thuộc đường tròn đường kính $BC$.
				\item $\triangle BFC$ vuông tại $F$ nên $B$, $F$, $C$ cùng thuộc đường tròn đường kính $BC$.
			\end{itemize}
			Suy ra $B$, $C$, $E$, $F$ cùng thuộc đường tròn đường kính $BC$.\\
			Vậy tứ giác $BEFC$ nội tiếp đường tròn.
			\item Ta có $\widehat{ACT}=90^{\circ}$ (góc nội tiếp chắn nửa đường tròn).\\
			Xét $\triangle ADB$ và $\triangle ACT$, có\\
			$\widehat{ACT}=\widehat{ADB}=90^{\circ}$.\\
			$\widehat{ABD}=\widehat{ATC}$ (góc nội tiếp cùng chắn $\stackrel\frown{AC}$).\\
			Vậy $\triangle ADB$ đồng dạng $\triangle ACT$ (g - g).\\
			Xét tứ giác $AEHF$, ta có
			\begin{itemize}
				\item $\triangle AEH$ vuông tại $E$ nên $A$, $E$, $H$ cùng thuộc đường tròn đường kính $AH$.
				\item $\triangle AFH$ vuông tại $F$ nên $A$, $F$, $H$ cùng thuộc đường tròn đường kính $AH$.
			\end{itemize}
			Suy ra $A$, $E$, $H$, $F$ cùng thuộc một đường tròn đường kính $AH$.
			Nên $\widehat{HEF}=\widehat{HAF}$ (góc nội tiếp cùng chắn $\stackrel\frown{HF}$).\\
			Mà $\widehat{DAB}=\widehat{CAT}$ (vì $\triangle ADB$ đồng dạng $\triangle ACT$) hay $\widehat{HAF}=\widehat{CAO}$.\\
			Suy ra $\widehat{HEF}=\widehat{CAO}$\quad ($1$).\\
			Lại có $OA=OC=R$ nên $\triangle CAO$ cân tại $O$.\\
			Khi đó $\widehat{CAO}=\widehat{ACO}$\quad ($2$).\\
			Xét $\triangle CAO$, có $\widehat{CAO}+\widehat{ACO}+\widehat{AOC}=180^{\circ}$\quad ($3$).\\
			Từ ($1$) và ($2$), ta được $\widehat{HEF}=\widehat{CAO}=\widehat{ACO}$.\quad ($4$)\\
			 Từ ($3$) và ($4$), ta được $\widehat{HEF}+\widehat{HEF}+\widehat{AOC}=180^{\circ}$ hay $2\widehat{HEF}+\widehat{AOC}=180^{\circ}$ (đpcm).
			\item Ta có $OB=OC=R$ nên $\triangle OBC$ cân tại $O$.\\
			Mà $\triangle OBC$ có $OM$ là đường trung tuyến suy ra $OM\perp BC$.\\
			Xét tứ giác $OMIC$, có
			\begin{itemize}
				\item $\triangle OMC$ vuông tại $M$ nên $O$, $M$, $I$ cùng thuộc đường tròn đường kính $OC$.
				\item $\triangle OIC$ vuông tại $I$ nên $O$, $I$, $C$ cùng thuộc đường tròn đường kính $OC$.
			\end{itemize}
			Suy ra bốn điểm $O$, $M$, $I$, $C$ cùng thuộc đường tròn đường kính $OC$.
			Nên $\widehat{CMI}=\widehat{COI}$ (góc nội tiếp cùng chắn $\stackrel\frown{CI}$).\\
			Ta có $\widehat{COI}$ là góc ở tâm chắn cung $\stackrel\frown{CT}$.\\
			$\widehat{CAT}$ là góc nội tiếp chắn cung $\stackrel\frown{CT}$.\\
			Suy ra $\widehat{COI}=2\widehat{CAT}$ (cùng chắn $\stackrel\frown{CT}$).\\
			Khi đó $\widehat{CMI}=2\widehat{DAB}$ (vì $\widehat{CAT}=\widehat{DAB}$).\quad ($1$)\\
			Ta có $\triangle BCF$ vuông tại $F$, có đường trung tuyến $FM$ nên $FM=MB$ suy ra $\triangle FMB$ cân tại $M$.\\
			Lại có $\widehat{FMB}=180^{\circ}-2\widehat{MBF}=2\left(90^{\circ}-\widehat{MBF}\right)=2\widehat{DAB}$\quad ($2$)\\
			Từ ($1$) và ($2$), ta được $\widehat{FMB}=\widehat{CMI}$.\\
			Mà $\widehat{FMB}+\widehat{FMC}=180^{\circ}$ nên $\widehat{CMI}+\widehat{FMC}=180^{\circ}$.\\
			Vậy $F$, $M$, $I$ thẳng hàng.
			\end{enumerate}
	}
\end{bt} 